% Generated draft: requires mathematical and source review.
\documentclass[11pt]{article}
\usepackage[margin=25mm]{geometry}
\usepackage{fontspec}
\setmainfont{Latin Modern Roman}
\usepackage{amsmath,amssymb,mathtools,unicode-math}
\setmathfont{Latin Modern Math}
\usepackage{enumitem,needspace,xurl,hyperref}
\setlength{\parindent}{0pt}
\setlength{\parskip}{5pt}
\setlength{\emergencystretch}{3em}
\newcommand{\N}{\mathbb{N}}
\newcommand{\Z}{\mathbb{Z}}
\newcommand{\Q}{\mathbb{Q}}
\newcommand{\R}{\mathbb{R}}
\newcommand{\C}{\mathbb{C}}
\newcommand{\F}{\mathbb{F}}
\newcommand{\A}{\mathbb{A}}
\newcommand{\PP}{\mathbb{P}}
\newcommand{\E}{\mathbb{E}}
\newcommand{\eps}{\varepsilon}
\newcommand{\dd}{\,\mathrm{d}}
\begin{document}
\section*{UnsolvedMath ID: 30003829}
\textbf{Source problem code:} catalogue entry
\subsection{Definitions and mathematical statement}
For a field \(F\), its Milnor \(K\)-theory is the graded ring
\[
K_*^M(F)=T_{\Z}(F^\times)\big/\langle a\otimes(1-a):a\in F\setminus\{0,1\}\rangle,
\]
where \(F^\times\) is regarded as an abelian group in degree one and \(T_{\Z}\) denotes its tensor algebra. This definition supplies terminology only: the source does not specify whether reconstruction uses this entire graded ring, selected degrees, a quotient, coefficients, or additional structure. Write \(\mathcal I(F)\) for the intended, unspecified Milnor \(K\)-theoretic reconstruction data.

A function field \(F/k\) means a finitely generated field extension; for an integral curve \(C\) of finite type over \(k\), its function field \(k(C)\) has transcendence degree one over \(k\). A perfect field is a field of characteristic zero or a field of positive characteristic whose Frobenius map is surjective. A \(k\)-variety here means an integral separated scheme of finite type over \(k\). It is normal if its local rings are integrally closed domains, proper if its structure morphism to \(\operatorname{Spec}k\) is proper, and geometrically integral if its base change to an algebraic closure of \(k\) is integral. Two integral \(k\)-varieties are birational over \(k\) if they have isomorphic dense open \(k\)-subvarieties.

\begin{enumerate}
\item \textbf{Curve counterexamples.} Can a counterexample to reconstruction from the intended data \(\mathcal I\) occur among function fields of curves over algebraically closed fields, or among function fields of curves over finite fields? A counterexample would have the schematic form
\[
\mathcal I(F)\cong\mathcal I(F'),\qquad F\not\cong F',
\]
with \(F\) and \(F'\) in the relevant curve function-field class. This is a question schema: the data, their permitted isomorphisms, whether the base field is fixed or identified, and whether the concluding field isomorphism must preserve the base field are not specified. The source also does not specify a geometric integrality or constant-field condition for these curves.

\item \textbf{Higher transcendence degree.} For which more general base fields \(k\), and under which additional hypotheses on a function field \(F/k\) with
\[
\operatorname{trdeg}_k F\geq 2,
\]
can Milnor \(K\)-theoretic reconstruction and its intended applications be extended? The supplied background describes proved finitely generated regular function-field cases, but does not state their exact hypotheses, conclusions, or reconstruction data. Here regularity of an extension means separability together with the condition that every element of \(F\) algebraic over \(k\) belongs to \(k\). The question does not specify whether regularity is to be retained in the proposed extensions, which additional base fields are intended, or which applications are to be extended.

\item \textbf{Reconstruction of varieties.} For a perfect field \(k\), what minimal data determine a normal, proper, geometrically integral \(k\)-variety \(X\) up to isomorphism, rather than merely up to birational equivalence? If the intended isomorphisms preserve a fixed base field, a proposed assignment of data \(\mathcal D\) would need to satisfy, for every pair \(X,Y\) in this class,
\[
\mathcal D(X)\cong\mathcal D(Y)\ \Longrightarrow\ X\cong_k Y.
\]
The source does not specify admissible assignments \(\mathcal D\), a comparison relation defining minimality, or whether these data must extend the preceding Milnor \(K\)-theoretic data. It also leaves the treatment of the base field implicit. Consequently this implication expresses a possible reconstruction criterion, not a fully specified conjecture about a particular invariant.
\end{enumerate}

These are three related research questions. The supplied record does not provide enough information to turn them into a unique, fully quantified reconstruction conjecture.
\subsection{Short English statement}
The source asks whether Milnor \(K\)-theoretic field reconstruction can fail for curve function fields over algebraically closed or finite fields, whether reconstruction extends to higher transcendence degree over broader base fields, and which minimal data recover a normal, proper, geometrically integral variety over a perfect field up to isomorphism. The reconstruction invariant, several hypotheses, and the meaning of minimality remain unspecified.
\subsection{Sources}
\begin{itemize}
\item Supplied UnsolvedMath local registry record, permanent ID \(30003829\), title \emph{Milnor K-Theory Reconstruction of Function Fields and Varieties}, exact source problem code \texttt{catalogue entry}, in the Oberwolfach Reports Open Problems collection. No individual proposer is supplied; the proposed year is recorded as 2018.
\item Supplied report citation: \emph{Field Arithmetic (2019)}, Oberwolfach Reports, DOI 10.4171/owr/2018/25, pp.\ 1517--1577. \url{https://doi.org/10.4171/owr/2018/25}. The supplied citation uses 2019, while the DOI and proposed-year metadata contain 2018; this dating has not been reconciled.
\item Supplied literature source labeled ``Cadoret--Pirutka'': \url{https://webusers.imj-prg.fr/~anna.cadoret/K2.pdf}. No title, theorem number, or precise reconstruction hypotheses are supplied.
\item Supplied literature source labeled ``Oberwolfach Report 2019'': \url{https://ems.press/content/serial-article-files/46747}. No precise location within the report is supplied for these questions.
\item Literature assessment dated exactly \texttt{2026-08-21}: the supplied curation labels the problem open and describes the broader reconstruction extensions and reconstruction up to isomorphism as unresolved beyond proved finitely generated regular function-field cases. This is supplied dated curation, not an independently verified literature assessment; the record itself requests independent verification of time-sensitive status and cited claims.
\item The supplied statement status is \texttt{corrected\_verified}; its accompanying explanation concerns preservation of mathematical content while repairing notation, typography, and sentence structure. That label does not establish a mathematical review of this draft or independent verification of the cited results.
\end{itemize}
\end{document}
